% GENERATED FILE. Edit canonical lessons, then run npm run generate:books.
% canonical-source-hash: fnv1a64:e1dd63332567ced4
% canonical-lessons: GU-C161-khasedvu, GU-C161-khankhervu, GU-C161-dhakvu, GU-C161-tutvu, GU-C161-bhegu-karvu

\chapter{Moving things about}
\label{ch:gu-moving-things-about}
\hlchaptermodality{161}

\begin{chapteropening}
\textbf{By the end of this chapter:} \emph{I can say what I move, cover and shake out at home.}

Complete the last lesson of chapter 161: I can say what I move, cover and shake out at home.
\end{chapteropening}

\section[\texorpdfstring{\gu{ખસેડવું} (khaseḍvũ)}{ખસેડવું (khaseḍvũ)}]{\gu{ખસેડવું} (khaseḍvũ) — to move, to shift}

\label{lesson:GU-C161-khasedvu}



\begin{quote}
Before the new one: say the Gujarati for electric light, then the Gujarati for to put out.
\end{quote}

\subsection*{You'll want to know: \gu{ખસેડવું}}
\textbf{\gu{ખસેડવું}} — \emph{khaseḍvũ} — \textquotedblleft{}to move, to shift\textquotedblright{}.

\begin{quote}
\textbf{\gu{અમે} \gu{પલંગ} \gu{ખસેડીએ} \gu{છીએ}.} — \emph{ame palaṅg khaseḍīe chīe} — \textquotedblleft{}We move the bed.\textquotedblright{}
\end{quote}

\begin{cousinweb}
\textbf{\gu{ખસેડવું}} (\emph{khaseḍvũ}) is built on the inherited \emph{khasvũ}, \textquotedblleft{}to slide, move\textquotedblright{}, with a causative ending. CDIAL traces \emph{khasvũ} to \emph{*khasati}, \textquotedblleft{}falls, slips\textquotedblright{}, a reconstructed form (CDIAL 3856).
\end{cousinweb}

\begin{grammarlens}[title={What you've built}]
One more word for everyday life at A2.
\end{grammarlens}

\subsection*{Your turn}
\begin{itemize}
\raggedright
  \item \emph{Say it:} \emph{khaseḍvũ}
  \item \emph{Say it:} \emph{khaseḍvũ}, once more
  \item \emph{Say it:} \emph{ame palaṅg khaseḍīe chīe} — We move the bed.
  \item \emph{Recall:} say the Gujarati for electric light, then the Gujarati for to put out, then say \emph{khaseḍvũ} again
\end{itemize}

\begin{tcolorbox}[breakable,colback=teal!4,colframe=teal!35!black,title={Before you move on}]
What does \gu{ખસેડવું} mean? (\textquotedblleft{}To move, to shift\textquotedblright{}.) Say it once more.
\end{tcolorbox}
\section[\texorpdfstring{\gu{ખંખેરવું} (khankhervũ)}{ખંખેરવું (khankhervũ)}]{\gu{ખંખેરવું} (khankhervũ) — to shake out, to dust off}

\label{lesson:GU-C161-khankhervu}



\begin{quote}
Before the new one: say the Gujarati for to put out, then the Gujarati for to move.
\end{quote}

\subsection*{You'll want to know: \gu{ખંખેરવું}}
\textbf{\gu{ખંખેરવું}} — \emph{khankhervũ} — \textquotedblleft{}to shake out, to dust off\textquotedblright{}.

\begin{quote}
\textbf{\gu{તે} \gu{ચાદર} \gu{ખંખેરે} \gu{છે}.} — \emph{te cādar khaṅkhere che} — \textquotedblleft{}She shakes the dust out of the sheet.\textquotedblright{}
\end{quote}

\begin{cousinweb}
\textbf{\gu{ખંખેરવું}} (\emph{khankhervũ}) is a tadbhava. CDIAL lists it, glossed \textquotedblleft{}to shake off dust\textquotedblright{}, under \emph{*khaṅkhālayati}, \textquotedblleft{}rinses\textquotedblright{}, a reconstructed form it takes to imitate a sound (CDIAL 3762).
\end{cousinweb}

\begin{grammarlens}[title={What you've built}]
One more word for everyday life at A2.
\end{grammarlens}

\subsection*{Your turn}
\begin{itemize}
\raggedright
  \item \emph{Say it:} \emph{khankhervũ}
  \item \emph{Say it:} \emph{khankhervũ}, once more
  \item \emph{Say it:} \emph{te cādar khaṅkhere che} — She shakes the dust out of the sheet.
  \item \emph{Recall:} say the Gujarati for to put out, then the Gujarati for to move, then say \emph{khankhervũ} again
\end{itemize}

\begin{tcolorbox}[breakable,colback=teal!4,colframe=teal!35!black,title={Before you move on}]
What does \gu{ખંખેરવું} mean? (\textquotedblleft{}To shake out, to dust off\textquotedblright{}.) Say it once more.
\end{tcolorbox}
\section[\texorpdfstring{\gu{ઢાંકવું} (ḍhā̃kvũ)}{ઢાંકવું (ḍhā̃kvũ)}]{\gu{ઢાંકવું} (ḍhā̃kvũ) — to cover}

\label{lesson:GU-C161-dhakvu}



\begin{quote}
Before the new one: say the Gujarati for to move, then the Gujarati for to shake out.
\end{quote}

\subsection*{You'll want to know: \gu{ઢાંકવું}}
\textbf{\gu{ઢાંકવું}} — \emph{ḍhā̃kvũ} — \textquotedblleft{}to cover\textquotedblright{}.

\begin{quote}
\textbf{\gu{માતા} \gu{તપેલી} \gu{ઢાંકે} \gu{છે}.} — \emph{mātā tapelī ḍhā̃ke che} — \textquotedblleft{}Mother covers the pot.\textquotedblright{}
\end{quote}

\begin{cousinweb}
\textbf{\gu{ઢાંકવું}} (\emph{ḍhā̃kvũ}) is a tadbhava. It comes from \emph{*ḍhakk}, \emph{*ḍhaṅk}, \textquotedblleft{}cover\textquotedblright{}, a reconstructed form beside Sanskrit \emph{ḍhakkana}, \textquotedblleft{}shutting\textquotedblright{} (CDIAL 5574).
\end{cousinweb}

\begin{grammarlens}[title={What you've built}]
One more word for everyday life at A2.
\end{grammarlens}

\subsection*{Your turn}
\begin{itemize}
\raggedright
  \item \emph{Say it:} \emph{ḍhā̃kvũ}
  \item \emph{Say it:} \emph{ḍhā̃kvũ}, once more
  \item \emph{Say it:} \emph{mātā tapelī ḍhā̃ke che} — Mother covers the pot.
  \item \emph{Recall:} say the Gujarati for to move, then the Gujarati for to shake out, then say \emph{ḍhā̃kvũ} again
\end{itemize}

\begin{tcolorbox}[breakable,colback=teal!4,colframe=teal!35!black,title={Before you move on}]
What does \gu{ઢાંકવું} mean? (\textquotedblleft{}To cover\textquotedblright{}.) Say it once more.
\end{tcolorbox}
\section[\texorpdfstring{\gu{તૂટવું} (tūṭvũ)}{તૂટવું (tūṭvũ)}]{\gu{તૂટવું} (tūṭvũ) — to break, to get broken}

\label{lesson:GU-C161-tutvu}



\begin{quote}
Before the new one: say the Gujarati for to shake out, then the Gujarati for to cover.
\end{quote}

\subsection*{You'll want to know: \gu{તૂટવું}}
\textbf{\gu{તૂટવું}} — \emph{tūṭvũ} — \textquotedblleft{}to break, to get broken\textquotedblright{}.

\begin{quote}
\textbf{\gu{પાતળી} \gu{સોય} \gu{જલદી} \gu{તૂટે} \gu{છે}.} — \emph{pātḷī soy jaldī tūṭe che} — \textquotedblleft{}A thin needle breaks quickly.\textquotedblright{}
\end{quote}

\begin{cousinweb}
\textbf{\gu{તૂટવું}} (\emph{tūṭvũ}) is a tadbhava. It comes from Sanskrit \emph{truṭyati}, \textquotedblleft{}is broken, falls apart\textquotedblright{}, and CDIAL lists the Gujarati \emph{tuṭvũ}, \textquotedblleft{}is broken\textquotedblright{} (CDIAL 6065).
\end{cousinweb}

\begin{grammarlens}[title={What you've built}]
One more word for everyday life at A2.
\end{grammarlens}

\subsection*{Your turn}
\begin{itemize}
\raggedright
  \item \emph{Say it:} \emph{tūṭvũ}
  \item \emph{Say it:} \emph{tūṭvũ}, once more
  \item \emph{Say it:} \emph{pātḷī soy jaldī tūṭe che} — A thin needle breaks quickly.
  \item \emph{Recall:} say the Gujarati for to shake out, then the Gujarati for to cover, then say \emph{tūṭvũ} again
\end{itemize}

\begin{tcolorbox}[breakable,colback=teal!4,colframe=teal!35!black,title={Before you move on}]
What does \gu{તૂટવું} mean? (\textquotedblleft{}To break, to get broken\textquotedblright{}.) Say it once more.
\end{tcolorbox}
\section[\texorpdfstring{\gu{ભેગું} \gu{કરવું} (bhegũ karvũ)}{ભેગું કરવું (bhegũ karvũ)}]{\gu{ભેગું} \gu{કરવું} (bhegũ karvũ) — to gather, to collect}

\label{lesson:GU-C161-bhegu-karvu}



\begin{quote}
Before the new one: say the Gujarati for to cover, then the Gujarati for to break.
\end{quote}

\subsection*{You'll want to know: \gu{ભેગું} \gu{કરવું}}
\textbf{\gu{ભેગું} \gu{કરવું}} — \emph{bhegũ karvũ} — \textquotedblleft{}to gather, to collect\textquotedblright{}.

\begin{quote}
\textbf{\gu{તે} \gu{કચરો} \gu{ભેગો} \gu{કરે} \gu{છે}.} — \emph{te kacro bhego kare che} — \textquotedblleft{}He gathers up the rubbish.\textquotedblright{}
\end{quote}

Here \textbf{\gu{ભેગો}} agrees with \textbf{\gu{કચરો}}.

\begin{cousinweb}
\textbf{\gu{ભેગું} \gu{કરવું}} (\emph{bhegũ karvũ}) is built from \emph{bhegũ}, \textquotedblleft{}together, collected\textquotedblright{}, and the inherited \emph{karvũ}, \textquotedblleft{}to do\textquotedblright{}. CDIAL lists \emph{bhegũ}, \textquotedblleft{}collected, mixed together\textquotedblright{}, under \emph{*bhēkkha}, \emph{*bhēgga}, \textquotedblleft{}lump\textquotedblright{}, a reconstructed form (CDIAL 9602).
\end{cousinweb}

\begin{grammarlens}[title={What you've built}]
One more word for everyday life at A2.
\end{grammarlens}

\subsection*{Your turn}
\begin{itemize}
\raggedright
  \item \emph{Say it:} \emph{bhegũ karvũ}
  \item \emph{Say it:} \emph{bhegũ karvũ}, once more
  \item \emph{Say it:} \emph{te kacro bhego kare che} — He gathers up the rubbish.
  \item \emph{Recall:} say the Gujarati for to cover, then the Gujarati for to break, then say \emph{bhegũ karvũ} again
\end{itemize}

\begin{tcolorbox}[breakable,colback=teal!4,colframe=teal!35!black,title={Before you move on}]
What does \gu{ભેગું} \gu{કરવું} mean? (\textquotedblleft{}To gather, to collect\textquotedblright{}.) Say it once more.
\end{tcolorbox}
